\section{Bitter Lesson 的正确理解}

讲师专门纠正了对 bitter lesson 的一种常见误读。
错误理解是：既然 scale 很重要，那算法就不重要了，只要一直加大资本开模型就行。
正确理解恰好相反：真正重要的是“能随着 scale 继续有效的算法”。

换句话说，最终模型效果可以粗略看成
\[
\text{accuracy} \approx \text{efficiency} \times \text{resources}.
\]

资源当然重要，但效率在大规模下更重要。
因为当训练成本已经到数千万、数亿美元时，任何浪费都会被放大。
讲师还引用了算法效率的历史改进，强调在很多任务上，算法进步本身可以带来数量级的收益，远远不只是硬件进步。

\begin{knowledgebox}{这门课的底层世界观}
\begin{itemize}
    \item 不是“算法无关”，而是“会随规模成长的算法最重要”。
    \item 研究的目标不是单纯堆资源，而是提高单位资源的产出。
    \item 问题可以被重新表述为：在给定 compute 和 data 预算下，能做出多好的模型？
\end{itemize}
\end{knowledgebox}

